% ======================================================================
% 09-local-weyl.tex (FULL REWRITE, Annals+++ / DIAMOND)
% Local Weyl law as a CANONICALITY + STABILITY statement (not just asymptotic)
% Locks inherited from Ch. 5–8:
%   - OBJ-LOCK: A = sqrt(Delta - 1/4), Tr^{reg} fixed by the Chapter-7 truncation ledger
%   - Normalization: T_\Gamma[h] = L_\Gamma[h - h(0)\chi] (Chapter 8)
%   - Error accounting: every admissible auxiliary change is absorbed into a finite ledger E_\Gamma
% ======================================================================

\chapter{A Local Weyl Law with Quantified Remainders}
\label{chap:09-local-weyl}

\epigraph{A Weyl law is not an asymptotic; it is a stability statement about what survives every admissible normalization.}{}

% ======================================================================
\section{Purpose and statement of the problem}
\label{sec:09.purpose}
% ======================================================================

\noindent\textbf{Scope.}
In the geometric setting fixed in Chapters~\ref{chap:01-introduction}--\ref{chap:08-normalization}, let
\[
A:=\sqrt{\Delta-\tfrac14}
\]
with the (regularized) trace functional $\Tr^{\mathrm{reg}}$ fixed by the truncation convention of Chapter~7.
This chapter extracts from the normalized trace functional of Chapter~8 a \emph{canonical principal term}
(the phase-volume constant) and a \emph{quantified remainder} that is stable under all admissible auxiliary choices.

\medskip
\noindent\textbf{Spectral parameterization.}
Write the discrete spectrum in spectral parameters $r_j\ge 0$ via
\[
\Delta\varphi_j=\Bigl(\tfrac14+r_j^2\Bigr)\varphi_j,
\qquad \{\varphi_j\}\ \text{orthonormal in }L^2.
\]
(When a continuous spectrum is present, it is handled inside $\Tr^{\mathrm{reg}}$ and the Chapter-7 cusp ledger.)

\medskip
\noindent\textbf{Spectral measures.}
We consider the unweighted and microlocally weighted measures
\begin{equation}\label{eq:09.measures}
\mu_\Gamma := \sum_j \delta_{r_j},
\qquad
\mu_\Gamma^{B} := \sum_j \langle B\varphi_j,\varphi_j\rangle\,\delta_{r_j},
\end{equation}
where $B\in\Psi^0(X)$ is a fixed zeroth-order pseudodifferential observable (Chapter~5).
The statement below shows that the leading term is \emph{canonical} (independent of auxiliary architecture)
and that all instabilities are absorbed into an explicit normalization ledger.

% ======================================================================
\section{Test functions and smoothing locks}
\label{sec:09.tests}
% ======================================================================

\subsection{Admissible test class}\label{subsec:09.admissible-smoothing}

Let $h:\R\to\C$ be admissible in the sense of Chapter~8 (Schwartz or Paley--Wiener under the locked Fourier conventions).
For $T\ge1$ define the rescaled family
\begin{equation}\label{eq:09.rescale}
h_T(t):=h(t/T).
\end{equation}

\medskip
\noindent\textbf{LOCK-F (Fourier convention).}
We use the Fourier transform
\begin{equation}\label{eq:09.FT}
\widehat h(\xi):=\int_{\R} h(t)\,e^{-it\xi}\,dt,
\end{equation}
consistent with Chapters~7--8.

\medskip
\noindent\textbf{Normalized trace object.}
Let $\chi$ be the admissible cutoff fixed in Chapter~8 ($\chi\equiv1$ near $0$). The canonical normalized trace is
\begin{equation}\label{eq:09.T.def}
\mathcal T_\Gamma[h;\chi]:=\mathcal L_\Gamma[h-h(0)\chi],
\end{equation}
with cutoff dependence controlled explicitly by Theorem~\ref{thm:08.normalization-stability}.

\subsection{Parity split and canonical parity reduction}\label{subsec:09.parity}

Define
\[
h^{\mathrm{ev}}(t)=\frac12(h(t)+h(-t)),
\qquad
h^{\mathrm{odd}}(t)=\frac12(h(t)-h(-t)),
\qquad h=h^{\mathrm{ev}}+h^{\mathrm{odd}}.
\]

\begin{lemma}[Parity reduction of the principal term]\label{lem:09.parity-reduction}
The principal Weyl term extracted from $\mathcal T_\Gamma[h;\chi]$ depends only on $h^{\mathrm{ev}}$.
Moreover, the odd-channel contribution carries no identity-sector mass and is therefore a pure remainder channel:
\[
\mathcal T_\Gamma[h;\chi]=\mathcal T_\Gamma[h^{\mathrm{ev}};\chi]+\mathcal T_\Gamma[h^{\mathrm{odd}};\chi],
\qquad
\mathcal T_\Gamma[h^{\mathrm{odd}};\chi]\ \text{has no identity contribution}.
\]
\end{lemma}

\begin{proof}
Linearity is immediate. The identity term in the local trace expansion (equivalently the local parametrix)
is an even functional of the test variable; hence it vanishes on $h^{\mathrm{odd}}$.
\end{proof}

% ======================================================================
\section{Smoothed spectral counting and local Weyl statements}
\label{sec:09.smoothed-weyl}
% ======================================================================

\subsection{Mollifiers and smoothed counts}\label{subsec:09.smoothed-count}

Fix $w\in\mathcal S(\R)$, $w\ge0$, with $\int_\R w(t)\,dt=1$. For $\delta\in(0,1]$ set
\[
w_\delta(t):=\delta^{-1}w(t/\delta).
\]
Define the smoothed counting functions
\begin{equation}\label{eq:09.N.smooth}
N_\Gamma(r;\delta):=(\mu_\Gamma*w_\delta)(r)=\sum_j w_\delta(r-r_j),
\end{equation}
and for $B\in\Psi^0(X)$,
\begin{equation}\label{eq:09.N.weighted}
N_{\Gamma,B}(r;\delta):=(\mu_\Gamma^{B}*w_\delta)(r)
=\sum_j \langle B\varphi_j,\varphi_j\rangle\,w_\delta(r-r_j).
\end{equation}

\subsection{Phase-volume constant (principal term)}\label{subsec:09.phase-volume}

Let $n=\dim X$ and let $d\omega$ denote Liouville measure on $S^*X$.
The universal phase-volume constant is
\begin{equation}\label{eq:09.CX.general}
C_X:=\frac{1}{(2\pi)^n}\,\vol(S^*X).
\end{equation}
In the surface case $n=2$ this simplifies to
\begin{equation}\label{eq:09.CX.surface}
C_X=\frac{\vol(X)}{2\pi}.
\end{equation}

\subsection{Smoothed local Weyl law (unweighted)}\label{subsec:09.statement}

\begin{theorem}[Smoothed local Weyl law with quantified remainders]\label{thm:09.smoothed-weyl}
Assume the standing hypotheses of Chapters~\ref{chap:01-introduction}--\ref{chap:08-normalization}.
Fix $w$ as above and $\delta\in(0,1]$.
Then as $r\to+\infty$,
\begin{equation}\label{eq:09.smoothed.weyl}
N_\Gamma(r;\delta)=C_X\,r^{n-1}+\mathcal R_\Gamma(r;\delta),
\end{equation}
where $C_X$ is given by \eqref{eq:09.CX.general} and the remainder admits the canonical two-part bound
\begin{equation}\label{eq:09.remainder.bound}
|\mathcal R_\Gamma(r;\delta)|
\le
\mathcal E^{\mathrm{TF}}_\Gamma(r;\delta)
+
\mathcal E^{\mathrm{Tab}}_\Gamma(r;\delta).
\end{equation}
Here:
\begin{enumerate}
\item $\mathcal E^{\mathrm{TF}}_\Gamma(r;\delta)$ is the trace-formula/normalization remainder controlled by the consolidated ledger
$\mathcal E_\Gamma[\cdot]$ of Chapter~8 (and by the Chapter-7 cusp ledger in the noncompact case);
\item $\mathcal E^{\mathrm{Tab}}_\Gamma(r;\delta)$ is the purely Tauberian smoothing error depending only on $w$ and $\delta$.
\end{enumerate}
In particular, the principal coefficient $C_X$ is canonical and independent of admissible auxiliary architecture.
\end{theorem}

\begin{remark}[Principal term as a locked invariant]\label{rem:09.principal-term}
The constant $C_X$ depends only on the principal symbol of $\Delta$ and the Riemannian volume measure; hence it is invariant under
admissible changes of quantization, partitions, and cutoffs. All such changes affect only $\mathcal E^{\mathrm{TF}}_\Gamma$.
\end{remark}

% ======================================================================
\section{Derivation from the normalized trace formula}
\label{sec:09.derivation}
% ======================================================================

\subsection{Spectral side under OBJ-LOCK}\label{subsec:09.spectral-side}

By Chapter~7 (spectral expansion with the fixed regularization),
\begin{equation}\label{eq:09.L.spectral}
\mathcal L_\Gamma[h]
=
\sum_j h(r_j)
+\mathcal C_\Gamma[h],
\end{equation}
where $\mathcal C_\Gamma[h]$ denotes the continuous-spectrum (and residual) contribution in the noncompact case
and is identically $0$ in the compact case.
Normalization yields
\[
\mathcal T_\Gamma[h;\chi]=\mathcal L_\Gamma[h-h(0)\chi],
\]
which removes the sole identity-sector ambiguity.

\subsection{Identity term and Plancherel density}\label{subsec:09.identity-term}

Let $\rho_X(r)$ denote the geometric Plancherel density (locked by convention). Then the identity contribution has the form
\begin{equation}\label{eq:09.identity.term}
\mathcal I_\Gamma[h]=\int_\R h(r)\,\rho_X(r)\,dr.
\end{equation}
For hyperbolic surfaces one may take $\rho_X(r)=\frac{\vol(X)}{2\pi}r\tanh(\pi r)$ (up to the Chapter-7 convention lock),
and the Weyl term is obtained from the large-$r$ expansion of $\rho_X(r)$.

\begin{lemma}[Identity contribution asymptotic for rescaled tests]\label{lem:09.identity-asymptotic}
Let $h_T$ be as in \eqref{eq:09.rescale}. Then, as $T\to+\infty$,
\begin{equation}\label{eq:09.identity.asymp}
\mathcal I_\Gamma[h_T]
=
C_X\,T^{n-1}\int_\R h(u)\,du
+
O\!\left(T^{n-2}\right),
\end{equation}
where $C_X$ is given by \eqref{eq:09.CX.general}.
\end{lemma}

\begin{proof}
This is the standard local Weyl expansion obtained from the wave kernel parametrix by Fourier inversion (stationary phase);
$C_X$ is determined by phase volume and is therefore canonical.
\end{proof}

\subsection{Non-identity terms and the trace-formula remainder ledger}\label{subsec:09.nonidentity}

Let $\mathcal G_\Gamma[h]$ denote the sum of non-identity geometric contributions (hyperbolic/orbit terms, elliptic terms, and
parabolic/scattering terms in the noncompact case) in the fixed Chapter-7 convention. Then
\begin{equation}\label{eq:09.T.decomp}
\mathcal T_\Gamma[h;\chi]=\mathcal I_\Gamma[h]+\mathcal G_\Gamma[h]+\mathcal R^{\mathrm{TF}}_\Gamma[h;\chi],
\end{equation}
where $\mathcal R^{\mathrm{TF}}_\Gamma[h;\chi]$ is the trace-formula remainder.

\begin{proposition}[Trace-formula remainder bound]\label{prop:09.trace-formula-remainder}
For admissible $h$,
\begin{equation}\label{eq:09.RTF.bound}
\big|\mathcal R^{\mathrm{TF}}_\Gamma[h;\chi]\big|
\le
\mathcal E_\Gamma[h],
\end{equation}
where $\mathcal E_\Gamma[h]$ is the consolidated normalization ledger of Definition~\ref{def:08.total-ledger}.
\end{proposition}

\begin{proof}
This is a direct application of Theorem~\ref{thm:08.normalization-stability} together with the fixed decomposition of
Chapter~7.
\end{proof}

% ======================================================================
\section{Two-component remainder decomposition: trace dust vs.\ Tauberian dust}
\label{sec:09.two-dust}
% ======================================================================

\subsection{Canonical dust split}\label{subsec:09.dust-split}

\begin{definition}[Two-component remainder decomposition]\label{def:09.two-component}
For $r>0$ and $\delta\in(0,1]$ define
\[
\mathcal R_\Gamma(r;\delta)=\mathcal R^{\mathrm{TF}}_\Gamma(r;\delta)+\mathcal R^{\mathrm{Tab}}_\Gamma(r;\delta),
\]
where $\mathcal R^{\mathrm{TF}}_\Gamma$ is the trace-formula/normalization remainder (geometry-dependent) and
$\mathcal R^{\mathrm{Tab}}_\Gamma$ is the Tauberian smoothing remainder (kernel-dependent).
\end{definition}

\subsection{Quantified Tauberian control}\label{subsec:09.tauberian}

\begin{lemma}[Tauberian remainder bound]\label{lem:09.tauberian-bound}
Let $w\ge0$ with $\int w=1$. Then for all $r>0$ and $\delta\in(0,1]$,
\begin{equation}\label{eq:09.tauberian.bound}
\big|N_\Gamma(r;\delta)-N_\Gamma(r)\big|
\le
N_\Gamma(r+\delta)-N_\Gamma(r-\delta).
\end{equation}
Consequently,
\[
\big|\mathcal R^{\mathrm{Tab}}_\Gamma(r;\delta)\big|
\le
N_\Gamma(r+\delta)-N_\Gamma(r-\delta).
\]
\end{lemma}

\begin{proof}
This is the standard monotonicity estimate for convolution with a nonnegative unit-mass kernel.
\end{proof}

% ======================================================================
\section{Weighted local Weyl laws (microlocal observables)}
\label{sec:09.weighted}
% ======================================================================

\subsection{Weighted counts}\label{subsec:09.observable}

Fix $B\in\Psi^0(X)$ with principal symbol $\sigma_0(B)$. Recall \eqref{eq:09.N.weighted}.

\begin{theorem}[Weighted smoothed local Weyl law]\label{thm:09.weighted-weyl}
Let $B\in\Psi^0(X)$ be fixed. Then
\begin{equation}\label{eq:09.weighted.weyl}
N_{\Gamma,B}(r;\delta)=C_{X,B}\,r^{n-1}+\mathcal R_{\Gamma,B}(r;\delta),
\qquad r\to+\infty,
\end{equation}
where
\begin{equation}\label{eq:09.CXB}
C_{X,B}:=\frac{1}{(2\pi)^n}\int_{S^*X}\sigma_0(B)\,d\omega,
\end{equation}
and
\begin{equation}\label{eq:09.weighted.remainder}
|\mathcal R_{\Gamma,B}(r;\delta)|
\le
\mathcal E^{\mathrm{TF}}_{\Gamma,B}(r;\delta)+\mathcal E^{\mathrm{Tab}}_{\Gamma,B}(r;\delta),
\end{equation}
with $\mathcal E^{\mathrm{TF}}_{\Gamma,B}$ controlled by the Chapter-8 normalization ledger (and $\mathcal E^{\mathrm{Tab}}_{\Gamma,B}$
by the Tauberian bound as in Lemma~\ref{lem:09.tauberian-bound}).
\end{theorem}

\begin{proof}
Insert $B$ into the microlocal trace expansion (Chapter~5) and apply the local trace/Weyl expansion to obtain the principal term
as the phase-space average of $\sigma_0(B)$. All quantization/partition dependence is lower order and is absorbed into the ledger;
the smoothing error is Tauberian.
\end{proof}

\begin{remark}[Quantization independence of $C_{X,B}$]\label{rem:09.quantization-independence}
The coefficient $C_{X,B}$ depends only on the principal symbol $\sigma_0(B)$; admissible changes of quantization alter only
subprincipal/lower-order contributions and thus affect only the remainder ledger.
\end{remark}

% ======================================================================
\section{Odd-channel defect estimate (intrinsic diagnostic)}
\label{sec:09.odd-defect}
% ======================================================================

\begin{proposition}[Odd-channel bound]\label{prop:09.odd-channel}
Let $h$ be admissible. Then
\begin{equation}\label{eq:09.odd.channel}
\mathcal T_\Gamma[h^{\mathrm{odd}};\chi]
=
\mathcal G_\Gamma[h^{\mathrm{odd}}]
+
\mathcal R^{\mathrm{TF}}_\Gamma[h^{\mathrm{odd}};\chi],
\end{equation}
and
\[
\big|\mathcal R^{\mathrm{TF}}_\Gamma[h^{\mathrm{odd}};\chi]\big|
\le
\mathcal E_\Gamma[h^{\mathrm{odd}}],
\]
so no identity-type Weyl main term occurs in the odd channel.
\end{proposition}

\begin{proof}
By Lemma~\ref{lem:09.parity-reduction}, the identity contribution vanishes on $h^{\mathrm{odd}}$.
The remainder estimate is Proposition~\ref{prop:09.trace-formula-remainder}.
\end{proof}

% ======================================================================
\section{From smoothed laws to sharp counting (Tauberian upgrade)}
\label{sec:09.sharp}
% ======================================================================

\begin{theorem}[Sharp Weyl law from uniform smoothed control]\label{thm:09.sharp-weyl}
Assume Theorem~\ref{thm:09.smoothed-weyl} holds uniformly for $\delta\in[r^{-\theta},1]$ for some fixed $\theta\in(0,1)$.
Then
\begin{equation}\label{eq:09.sharp.weyl}
N_\Gamma(r)=C_X\,r^{n-1}+o(r^{n-1}),
\qquad r\to+\infty.
\end{equation}
More precisely, for any admissible choice $\delta(r)\ge r^{-\theta}$,
\begin{equation}\label{eq:09.sharp.bound}
N_\Gamma(r)
=
C_X\,r^{n-1}
+
O\!\big(\mathcal E^{\mathrm{TF}}_\Gamma(r;\delta(r))\big)
+
O\!\big(N_\Gamma(r+\delta(r))-N_\Gamma(r-\delta(r))\big).
\end{equation}
\end{theorem}

\begin{proof}
Combine Theorem~\ref{thm:09.smoothed-weyl} with Lemma~\ref{lem:09.tauberian-bound}.
Choose $\delta(r)$ so that the boundary-layer increment is negligible compared with $r^{n-1}$; the trace remainder is uniformly
controlled by the normalization ledger.
\end{proof}

% ======================================================================
\section{Uniformity in auxiliary data (portability clause)}
\label{sec:09.uniformity}
% ======================================================================

\begin{corollary}[Stability under admissible auxiliary changes]\label{cor:09.stability}
Let $\mathcal T_\Gamma[h;\chi]$ and $\mathcal T'_\Gamma[h;\chi']$ be normalized trace functionals built from two admissible
auxiliary architectures. Then the Weyl coefficients $C_X$ and $C_{X,B}$ in
Theorems~\ref{thm:09.smoothed-weyl} and~\ref{thm:09.weighted-weyl} coincide, and the difference between the corresponding remainders
is bounded by the sum of the associated consolidated ledgers (Chapter~8), modulo the explicit identity-sector correction.
\end{corollary}

\begin{proof}
The principal coefficients arise from the identity term (phase volume) and are therefore canonical. Remainder stability follows
from Theorem~\ref{thm:08.normalization-stability}.
\end{proof}

% ======================================================================
\section{Bibliographic notes}\label{sec:09.bib-notes}
% ======================================================================

Classical and microlocal Weyl laws are standard; see \cite{HormanderI,Sogge2014,Zworski2012}.
In automorphic/hyperbolic settings, the spectral parameterization and principal term are consistent with the Selberg trace formula;
see \cite{Hejhal1983,Buser1992,Iwaniec2002,Borthwick2016}.
Weighted Weyl laws for pseudodifferential observables follow from symbolic calculus and local trace expansions; see
\cite{DuistermaatGuillemin1975,Ivrii1998,DyatlovZworski2019}.

% ======================================================================
% Audit (Chapter 9, Weyl layer)
% ======================================================================

\section*{Audit (Local Weyl layer)}\label{sec:09.audit}

\noindent\textbf{Locks.}
\begin{itemize}
\item[(L9.1)] OBJ-LOCK: $A=\sqrt{\Delta-\tfrac14}$ and $\Tr^{\mathrm{reg}}$ fixed by Chapter~7.
\item[(L9.2)] Normalization lock: $\mathcal T_\Gamma[h;\chi]=\mathcal L_\Gamma[h-h(0)\chi]$ (Chapter~8).
\item[(L9.3)] Principal invariant: $C_X=(2\pi)^{-n}\vol(S^*X)$, hence $C_X=\vol(X)/(2\pi)$ for $n=2$.
\item[(L9.4)] Remainder split: $\mathcal R=\mathcal R^{\mathrm{TF}}+\mathcal R^{\mathrm{Tab}}$ (geometry vs smoothing).
\item[(L9.5)] Odd-channel diagnostic: no identity/Weyl term on $h^{\mathrm{odd}}$.
\end{itemize}

\medskip
\noindent\textbf{Failure modes.}
\begin{itemize}
\item[(F9.1)] Treating $\mathcal E_\Gamma$ as optional $\Rightarrow$ false portability (auxiliary dependence leaks).
\item[(F9.2)] Mixing Fourier/sign conventions between Chapters~7 and~9 $\Rightarrow$ wrong $C_X$ and wrong scaling.
\item[(F9.3)] Confusing Tauberian dust with trace dust $\Rightarrow$ incorrect remainder claims.
\end{itemize}

\medskip
\noindent\textbf{STATUS\_CERT.}
CLOSED (structural): canonical principal term and quantified remainder structure fixed; uniformity under auxiliary changes reduced to
Chapter-8 ledger plus explicit identity correction.
NEXT(1): specialize $\mathcal E^{\mathrm{TF}}_\Gamma(r;\delta)$ to the window family $h_{\lambda,\eta}$ (with $\delta\sim\eta$)
and record the explicit $(\lambda,\eta)$-scaling needed for the application chapters.

% ======================================================================
% End of 09-local-weyl.tex
% ======================================================================
